\documentclass[a4paper,11pt]{article}

% Always keep your main language last
\usepackage[british,nynorsk,samin,norsk]{babel}

\usepackage{_sty/UiT}
\usepackage{_sty/IMS}
\usepackage{MAT-0001/MAT-0001}

%==============================================================================%
%                              COURSE INFORMATION                              %
%==============================================================================%

\settoggle{isLF}{true}
% \settoggle{isLF}{false} % Comment this line to show solutions

\frontpageUiTsetup{%
  Language = {norsk}, %auto = language defined by babel
  isExam = {false}, % Is this an exam? true / false
  isExercise = {true}, % Is this an exercise? true / false
  ExerciseNumber = {1}, % Which exercise number is this?
  CourseCode = {auto}, % auto = \courseCode (defined in courseCode.sty, e.g. MAT-1001.sty)
  CourseName = {auto}, % auto = \courseName (defined in courseCode.sty, e.g. MAT-1001.sty)
  Year = {2019}, % No auto setting for this, as it might change on recompilation
  Month = {8}, % No auto setting for this, as it might change on recompilation
  Day = {27},% No auto setting for this, as it might change on recompilation
  %    DeadlineYear = {2020},
  %    DeadlineMonth = {12},
  %    DeadlineDay = {22},
  DurationDays = {3}, % Set the number of days the exam should last
}

%==============================================================================%
%                     OWN COMMANDS AND PACKAGES BELOW HERE                     %
%==============================================================================%

\def\firstcircle{(0,0) circle(1cm)}
\def\secondcircle{(0:1cm) circle (1cm)}
\def\middlecircle{(0:1cm) circle (0.5cm)}

\excludecomment{solution}

\begin{document}

%==============================================================================%
%                                 EXERCISE 01                                  %
%==============================================================================%

\frontpageUiT

\settoggle{isLF}{false}
\settoggle{isLFshort}{false}
\excludecomment{solution}

\titlebox[norsk]{auto}{\Spraak}{auto}{\dateDuration}

%==============================================================================%
%                                 SEKSJON 1.1                                  %
%==============================================================================%

\subsection*{Seksjon 1.1}

\begin{problem}[1]
  La $A = \Set{3, 4, 5}$ og $B = \Set{3, 5}$. Hvilke 
  utsagn er sanne?
  %
  \begin{tabenum}[\color{\UiTnumbercolor}\bfseries a\color{black})]%
    \tabenumitem {$3 \in B$} \UiTcorrect
    \tabenumitem {$4, 5 \in A$} \UiTcorrect
    \tabenumitem {$2 \in A$} \UiTwrong \\[1ex]
    \tabenumitem {$B \subseteq A$} \UiTcorrect
    \tabenumitem {$A \subseteq B$} \UiTwrong
    \tabenumitem {$7 \not\in A$} \UiTcorrect
  \end{tabenum}
  %
\end{problem}

\excludecomment{solution}

\begin{solution}
  Svarene står ovenfor, men la oss raskt forklare symbolene. $\in$ betyr
  \bquote{element}, og $3 \in B$ spør enkelt og greit om $3$ er et element i
  $B$.  Med andre ord om $3$ er inneholdt i $B$.

  Symbolet $\subseteq$ minner om krokodillemunnen fra barneskolen, og betyr
  \emph{delmengde} eller \emph{undermengde}. $B \subseteq A$ spør om \emph{alle}
  elementene i $A$ er inneholdt i $B$, dersom dette er sant er $A$ en
  undermengde av $B$.
\end{solution}

\begin{problem}
  Gitt mengden $M = \Set{1, 2, 3, 4, 5, 6, 7}$. Skriv følgende
  mengder på listeform:
    \begin{subproblem}{2}
      \item $\Set{x \in M \given x \neq 2 \ \text{og} \ x \neq 5}$
        \label{subproblem:1.1-2a-oving-01-2019-MAT-0001}
      \item $\Set{x + 7 \given x \in M}$
        \label{subproblem:1.1-2b-oving-01-2019-MAT-0001}
    \end{subproblem}
\end{problem}

\begin{solution}
  \Opg{subproblem:1.1-2a-oving-01-2019-MAT-0001} Ønsker å skrive mengden
  $\Set{x \in M \given x \neq 2 \ \text{og} \ x \neq 5}$ på listeform.
  Notasjonen kan og leses slik: Ønsker alle elementer i $M$ slik at elementene
  er ulik $2$ og $5$. Hvor tegnet $\mid$ leses som \bquote{slik at}.  Dette er
  selvsagt mengden
  %
  \begin{equation*}
    \Set{1, 3, 4, 6, 7}
  \end{equation*}
  %
  \Opg{subproblem:1.1-2b-oving-01-2019-MAT-0001} Ønsker å skrive mengden
  $\Set{x + 7 \given x \in M}$ på listeform. Her kan vi tenke oss at vi ønsker å
  gjøre $x + 7$ gitt at $x$ ligger i $M$.  Altså må vi ta alle verdiene i
  mengden $M$ og legge til $7$:
  %
  \begin{equation*}
    \Set{8, 9, 10, 11, 12, 13, 14}
  \end{equation*}
\end{solution}

\newpageLF

\begin{problem}
  La $R = \Set{1, 3, 5, 7}$ og $S = \Set{3, 5, 7, 9}$.
  Finn $R \cap S$, $R \cup S$, $R \setminus S$, $S \setminus R$.
\end{problem}

\begin{solution}
    \def\firstcircle{(0,0) circle(1cm)}
    \def\secondcircle{(0:1cm) circle (1cm)}
    %
  \textcolor{UiT-main}{$R \cup S$}: Husk at du kan tenke på $\cup$ som addisjon.
  Vi ønsker å legge sammen alle elementene i $R$ og $S$. \textbf{NB:} husk at i
  en mengde er alle elementene unike så vi skriver ikke ett element mer enn èn
  gang.

  \begin{minipage}{0.5\textwidth}
    \centering
    \begin{align*}
      R & = \Set{1, 3, 5, 7} \\
      S & = \quad \Set{3, 5, 7, 9} \\
      R \cup S & = \Set{1, 3, 5, 7, 9}
    \end{align*}
  \end{minipage}
  \begin{minipage}{0.5\textwidth}
    \centering
    \begin{tikzpicture}
        \fill[opacity=0.75,UiT-blue] \firstcircle;
        \fill[opacity=0.75,UiT-blue] \secondcircle;
        \draw[thick] \firstcircle node[left] {\color{black}$R$};
        \draw[thick] \secondcircle node[right]{\color{black}$S$};
        \node at (0.5, 0)  (m)    {$\cup$};
        \pgfresetboundingbox
        \path[use as bounding box] (-1,-1) rectangle (2,1);
    \end{tikzpicture}
  \end{minipage} \medskip

  \textcolor{UiT-main}{$R \cap S$}: Her leses symbolet $\cap$ som \emph{snitt},
  det er altså snakk bare det som ligger i \emph{begge} mengdene. Dersom noe
  ligger i bare en av mengdene skal det altså ikke tas med.

  \begin{minipage}{0.5\textwidth}
    \centering
    \begin{align*}
      R & = \Set{1, 3, 5, 7} \\
      S & = \quad \Set{3, 5, 7, 9} \\
      R \cap S & = \quad \Set{3, 5, 7}
    \end{align*}
  \end{minipage}
  \begin{minipage}{0.5\textwidth}
    \centering
    \begin{tikzpicture}
        \begin{scope}
          \clip \firstcircle;
          \fill[opacity=0.75,fill=UiT-blue] \secondcircle;
        \end{scope}
        \draw[thick] \firstcircle node[left] {\color{black}$R$};
        \draw[thick] \secondcircle node[right]{\color{black}$S$};
        \node at (0.5, 0)  (m)    {$\cap$};
        \pgfresetboundingbox
        \path[use as bounding box] (-1,-1) rectangle (2,1);
    \end{tikzpicture}
  \end{minipage} \medskip

  \textcolor{UiT-main}{$R \setminus S$}: Her leses $\setminus$ som
  \emph{komplimentet} til.  Dette kan en litt feilaktig tenke på som
  \emph{subtraksjon} eller minus. Vi ønsker med andre ord alt som ligger i $R$
  men \emph{ikke} i $S$:

  \begin{minipage}{0.5\textwidth}
    \centering
    \begin{align*}
      R & = \Set{1, 3, 5, 7} \\
      S & =   \quad \Set{3, 5, 7, 9} \\
      R \setminus S & =  \Set{1}
    \end{align*}
  \end{minipage}
  \begin{minipage}{0.5\textwidth}
    \centering
    \begin{tikzpicture}
        \begin{scope}[even odd rule]% first circle without the second
            \clip \secondcircle (-1.5,-1.6) rectangle (2,1.6);
            \fill[UiT-blue,opacity=0.75] \firstcircle;
        \end{scope}
        \draw[thick] \firstcircle node[left,xshift=0.35em] {\color{black}$R\setminus S$};
        \draw[thick] \secondcircle node[right]{\color{black}$S$};
        \pgfresetboundingbox
        \path[use as bounding box] (-1,-1) rectangle (2,1);
    \end{tikzpicture}
  \end{minipage} \medskip

  \textcolor{UiT-main}{$S \setminus R$}: Akkuratt som i forrige deloppgave
  ønsker vi nå alt som ligger i mengden $S$, men ikke i $R$

  \begin{minipage}{0.5\textwidth}
    \centering
    \begin{align*}
      R & = \Set{1, 3, 5, 7} \\
      S & =   \quad \Set{3, 5, 7, 9} \\
      S \setminus R & =  \qquad \qquad \Set{9}
    \end{align*}
  \end{minipage}
  \begin{minipage}{0.5\textwidth}
    \centering
    \begin{tikzpicture}
        \begin{scope}[even odd rule]% first circle without the second
            \clip \firstcircle (-1,-1.6) rectangle (3.5,1.6);
            \fill[UiT-blue,opacity=0.75] \secondcircle;
        \end{scope}
        \draw[thick] \firstcircle node[left] {\color{black}$R$};
        \draw[thick] \secondcircle node[right,xshift=-0.35em] {\color{black}$S\setminus R$};
        \pgfresetboundingbox
        \path[use as bounding box] (-1,-1) rectangle (2,1);
    \end{tikzpicture}
  \end{minipage} \medskip
\end{solution}

\newpageLF

%==============================================================================%
%                                 SEKSJON 1.2                                  %
%==============================================================================%

\subsection*{Seksjon 1.2}

%==============================================================================%
%                                PROBLEM 1.2.1                                 %
%==============================================================================%

\begin{problem}[1]
  Skriv følgende mengder som intervaller:
  %
  \begin{subproblem}{2}
    \item $\Set{x \in \R \given        x \leq 10}$
      \label{subproblem:1.2-1a-oving-01-2019-MAT-0001}
    \item $\Set{x \in \R \given 2 \leq x \leq  5}$
      \label{subproblem:1.2-1b-oving-01-2019-MAT-0001}
  \end{subproblem}
\end{problem}

\begin{solution}
  \Opg{subproblem:1.2-1a-oving-01-2019-MAT-0001}
  Her ønsker vi enkelt og greit alle reelle tall mindre eller lik $10$:
  %
  \begin{equation*}
    \Set{x \in \R \given        x \leq 10}
    \quad \Leftrightarrow \quad
    x \in (\infty,10],
  \end{equation*}
  %
  hvor $($ ble brukt for å \emph{ekskludere} $\infty$\footnote{Grunnen til at vi
  må ekskludere $\infty$ er at dette ikke er ett tall, men en grenseverdi vi
bare kan \emph{nærme} oss, men aldri nå.  Det er derfor vi bruker opp til, men
ikke inkludert.} mens $]$ ble brukt for å \emph{inkludere} $10$.
\end{solution}

\medskip

\begin{solution}
  \Opg{subproblem:1.2-1b-oving-01-2019-MAT-0001}
  Mengden $\Set{x \in \R \given 2 \leq x \leq  5}$ kan leses som at vi ønsker
  alle relle tall som er større eller lik $2$ og mindre eller lik $5$:
  %
  \begin{equation*}
    \Set{x \in \R \given 2 \leq x \leq  5}
    \quad \Leftrightarrow \quad
    x \in [2,5]
  \end{equation*}
  %
  hvor $[$ ble brukt for å vise at dette er ett \emph{lukket} intervall; 
  altså ett intervall der endepunktet er inkludert. 
\end{solution}

%==============================================================================%
%                                PROBLEM 1.2.2                                 %
%==============================================================================%

\begin{problem}
  Hvilke mengder er like?
  %
  \begin{alignat*}{2}
    A &= \Set{x\in\R\given2\leq x<5} \qquad && E=(2,\infty)\\
    B &= \Set{x\in\R\given2 < x\leq5}     && F=[2,5)\\
    C &= \Set{x\in\R\given x\leq5}          && G=(2,5]\\
    D &= \Set{x\in\R\given2<x}              && H=(-\infty,5]
  \end{alignat*}
\end{problem}

\begin{solution}
  Det enkleste er nok enten å skrive mengdene som intervaller, eller skrive alle
  intervallene som mengder. Dette er akkurat det samme, men vi velger å gjøre
  førstnevnte:
  %
  \begin{alignat*}{2}
    A &= \Set{x\in\R\given2\leq x<5} \quad && \Leftrightarrow \quad x \in [2,5) = F\\
    B &= \Set{x\in\R\given2< x\leq5}       && \Leftrightarrow \quad x \in (2,5] = G \\
    C &= \Set{x\in\R\given x\leq5}         && \Leftrightarrow \quad x \in (-\infty,5] = H\\
    D &= \Set{x\in\R\given2<x}             && \Leftrightarrow \quad x \in (2,\infty) = E
  \end{alignat*}
  %
  Her byttet vi i praksis ut hver $\leq$ med enten $[$ eller $]$, og $<$ med
  enten $($ eller $)$.
\end{solution}

%==============================================================================%
%                                PROBLEM 1.2.3                                 %
%==============================================================================%

\begin{problem}
  \label{problem:1.2-1b-oving-01-2019-MAT-0001} Skriv følgende uttrykk som et
  enkelt intervall (eller som den tomme mengden, $\varnothing$, om nødvendig)
  %
  \begin{tabenum}[\color{\UiTnumbercolor}\bfseries a\color{black})]%
    \tabenumitem \hypertarget{subproblem:1.2-1b-oving-01-2019-MAT-0001}{$[1, 3] \cap (2, 4]$}
    \tabenumitem \hypertarget{subproblem:1.2-1b-oving-01-2019-MAT-0001}{$[1, 3] \cup (2, 4]$} \\[1ex]
    \tabenumitem \hypertarget{subproblem:1.2-1b-oving-01-2019-MAT-0001}{$[1, 3] \cap [3, 4]$}
    \tabenumitem \hypertarget{subproblem:1.2-1b-oving-01-2019-MAT-0001}{$(1, 3) \cap (3, 4)$} \\[1ex]
    \tabenumitem \hypertarget{subproblem:1.2-1b-oving-01-2019-MAT-0001}{$(-\infty, 3) \cap (2, \infty)$}
    \tabenumitem \hypertarget{subproblem:1.2-1b-oving-01-2019-MAT-0001}{$(-\infty,30) \cap \Set{30}$}
  \end{tabenum}
\end{problem}

\hvisLF{%
  \begin{figure*}[htbp!]
    \centering
    \begin{subfigure}[t]{0.5\textwidth}
      \centering
      \begin{tikzpicture}[scale=1]
        \standardAxis
        \drawIntervalCC{1}{3}{UiT-blue}
        \drawIntervalOC{2}{4}{UiT-orange}
        \fill[opacity = 0.2, green,rounded corners=1ex]%
        (2,-1.12ex) -- (3, -1.12ex) -- (3, 1.12ex) -- (2,1.12ex) -- cycle;
      \end{tikzpicture}
      \caption{$\textcolor{UiT-blue}{[1, 3]} \cap \textcolor{UiT-orange}{(2,4]} =  \UiTcorrectColor{(2,3]}$}
    \end{subfigure}%
    ~ 
    \begin{subfigure}[t]{0.5\textwidth}
      \centering
      \begin{tikzpicture}[scale=1]
        \standardAxis
        \drawIntervalCC{1}{3}{UiT-blue} 
        \drawIntervalCC{2}{4}{UiT-blue}
        \fill[opacity = 0.2, green,rounded corners=1ex]%
        (1,-1.12ex) -- (4, -1.12ex) -- (4, 1.12ex) -- (1,1.12ex) -- cycle;
      \end{tikzpicture}
      \caption{$\textcolor{UiT-blue}{[1, 3]} \cup \textcolor{UiT-orange}{(2,4]} =  \UiTcorrectColor{[1,4]}$}
    \end{subfigure}

    \vspace{\baselineskip}

    \begin{subfigure}[t]{0.5\textwidth}
      \centering
      \begin{tikzpicture}[scale=1]
        \standardAxis
        \drawIntervalCC{1}{3}{UiT-blue} 
        \drawIntervalCC{3}{4}{UiT-orange}
      \end{tikzpicture}
      \caption{$\textcolor{UiT-blue}{[1, 3]} \cap \textcolor{UiT-orange}{[3,4]} =  \UiTcorrectColor{[3,3]}$}
      \label{fig:1.2-3c-MAT-0001-2019-V}
    \end{subfigure}%
    ~ 
    \begin{subfigure}[t]{0.5\textwidth}
      \centering
      \begin{tikzpicture}[scale=1]
        \standardAxis
        \drawIntervalOO{1}{3}{UiT-blue} 
        \drawIntervalOO{3}{4}{UiT-orange}
      \end{tikzpicture}
      \caption{$\textcolor{UiT-blue}{(1, 3)} \cap \textcolor{UiT-orange}{(3,4)} =  \UiTcorrectColor{\varnothing}$}
    \end{subfigure}

    \vspace{\baselineskip}

    \begin{subfigure}[t]{0.5\textwidth}
      \centering
      \begin{tikzpicture}[scale=1]
        \standardAxis
        \draw[{-)}, ultra thick, blue] (2.99,0) -- (3,0);
        \fill[opacity = 0.2, UiT-blue,rounded corners=1ex]%
        (0.5,-1.12ex) -- (3, -1.12ex) -- (3, 1.12ex) -- (0.5,1.12ex) -- cycle;
        \draw[{(-}, ultra thick, UiT-orange] (2,0) -- (2.01,0);
        \fill[opacity = 0.2, UiT-orange,rounded corners=1ex]%
        (2,-1.12ex) -- (4.5, -1.12ex) -- (4.5, 1.12ex) -- (2,1.12ex) -- cycle;
        \fill[opacity = 0.2, green,rounded corners=1ex]%
        (2,-1.12ex) -- (3, -1.12ex) -- (3, 1.12ex) -- (2,1.12ex) -- cycle;
      \end{tikzpicture}
      \caption{$\textcolor{UiT-blue}{(-\infty,3)} \cap \textcolor{UiT-orange}{(2, \infty)} = \UiTcorrectColor{(2,3)}$}
    \end{subfigure}%
    ~ 
    \begin{subfigure}[t]{0.5\textwidth}
      \centering
      \begin{tikzpicture}[scale=1]
        \draw[->, thick] (27.5,0) -- (31.5,0);
        \foreach \x/\xtext in {28/28,29/29,30/30,31/31}
        \draw[thick] (\x,1pt) -- (\x,-1pt) node[below] {\xtext};
        \pgfmathsetmacro\asnub{0.999*30}
        \draw[{-)}, ultra thick, UiT-blue] (\asnub,0) -- (30,0);
        \fill[opacity = 0.2, UiT-blue,rounded corners=1ex]%
        (27.5,-1.12ex) -- (30, -1.12ex) -- (30, 1.12ex) -- (27.5,1.12ex) -- cycle;
        \draw[{(-}, ultra thick, UiT-orange] (29.9,0) -- (30,0);
      \end{tikzpicture}
      \caption{$\textcolor{UiT-blue}{(-\infty,30)} \cap \textcolor{UiT-orange}{\{30\}} =  \UiTcorrectColor{\varnothing}$}
    \end{subfigure}
    \caption{Illustrasjoner av intervallene fra \cref{problem:1.2-1b-oving-01-2019-MAT-0001}.}
    \label{fig:1.2-3-MAT-0001-2019-V}
  \end{figure*}
}

\begin{solution}
  \Opg{problem:1.2-1b-oving-01-2019-MAT-0001}: Her kan det lønne seg
  og tegne intervallene slik som vi har gjort i
  \cref{fig:1.2-3-MAT-0001-2019-V}. Om en ikke har en Michelangelo\footnote{Her
    er det altså snakk om kunstneren ikke en av de fire hovedkarakteren fra den
  svært engasjerende serien Teenage Mutant Ninja Turtles.} i magen, kan en også
  skrive opp intervallene som ulikheter, og sortere dem:
  %
  \begin{flalign*}
    \hyperlink{subproblem:1.2-1a-oving-01-2019-MAT-0001}{\textcolor{\UiTrefcolor}{a})}
        && [1, 3] \cap (2, 4] 
        \quad \Leftrightarrow \quad
    1 \UiTcorrectColor{< 2 \leq 3}{} \leq 4 &&
  \end{flalign*}
  %
  Siden det er snakk om \emph{snittet}, ønsker vi det minste intervalet, altså
  $(2,3]$. 
  %
  \begin{flalign*}
    \hyperlink{subproblem:1.2-1b-oving-01-2019-MAT-0001}{\textcolor{\UiTrefcolor}{b})}
        && [1, 3] \cap (2, 4] 
        \quad \Leftrightarrow \quad
    \UiTcorrectColor{1 < 2 \leq 3 \leq 4}{} &&
  \end{flalign*}
  %
  Siden det er snakk om \emph{union} ønsker vi det \emph{største} snittet, altså
  $[1,4]$. 
  %
  \begin{flalign*}
    \hyperlink{subproblem:1.2-1c-oving-01-2019-MAT-0001}{\textcolor{\UiTrefcolor}{c})}
        && [1, 3] \cap [3, 4] 
        \quad \Leftrightarrow \quad
    1 \UiTcorrectColor{\leq 3 \leq 3}{} \leq 4 &&
  \end{flalign*}
  %
  Her får vi nå ett intervall der vi skal ha alle tall fra og med $3$ til og med
  $3$, det eneste intervallet som tilfredstiller dette er selvsagt $[3,3]$. For
  å være veldig pedantisk blir det derimot feil å si at intervallet er $\{3\}$.
  Intervallet er $[3,3]$ som bare inneholder ett punkt, nemlig $\{3\}$. Som en
  og ser fra \cref{fig:1.2-3c-MAT-0001-2019-V}.
  %
  \begin{flalign*}
    \hyperlink{subproblem:1.2-1d-oving-01-2019-MAT-0001}{\textcolor{\UiTrefcolor}{d})}
        && (1, 3) \cap (3, 4) 
        \quad \Leftrightarrow \quad
    1 \UiTcorrectColor{< 3 < 3}{} < 4 &&
  \end{flalign*}
  %
  Her ønsker vi å finne intervallet som inneholder alle tall \emph{større} enn
  $3$, men som også er \emph{mindre} enn $3$. Ah, ha! En umulighet! Siden ingen
  tall kan oppfylle dette kravet skriver vi det som den \emph{tomme mengden}
  $(1, 3) \cap (3, 4) =\varnothing$.
  %
  \begin{flalign*}
    \hyperlink{subproblem:1.2-1e-oving-01-2019-MAT-0001}{\textcolor{\UiTrefcolor}{e})}
        && (-\infty, 3) \cap [2, \infty]
        \quad \Leftrightarrow \quad
    - \infty \UiTcorrectColor{< 2 \leq 3 <} \infty &&
  \end{flalign*}
  %
  Her må vi passe på å bruke $<$ når det kommer til $\infty$ for dette er en
  grense, ikke ett tall. Dette gir da intervallet $(2,3)$.
  %
  \begin{flalign*}
    \hyperlink{subproblem:1.2-1f-oving-01-2019-MAT-0001}{\textcolor{\UiTrefcolor}{f})}
        && (-\infty, 30) \cap \Set{30}
        \quad \Leftrightarrow \quad
    - \infty \UiTcorrectColor{< 30 < 30} &&
  \end{flalign*}
  %
  Hvor vi da ser at det minste intervallet er den tomme mengden
  \UiTcorrectColor{$\varnothing$}.  Dette er fordi intervallet vårt går opp til men
  \emph{ikke} inkludert $30$.
\end{solution}


%==============================================================================%
%                                PROBLEM 1.2.4                                 %
%==============================================================================%

\begin{problem}
  \label{problem:1.2-4-oving-01-2019-MAT-0001} I teksten foran ble det gitt en
  formell definisjon av begrepet intervall. Bruk denne definisjonen til å bevise
  at dersom $U$ og $V$ er to intervaller som har minst to tall til felles, så er
  også snittet $U \cap V$ et intervall.
\end{problem}

\hvisLF{%
  \begin{figure*}[htbp!]
    \centering
    \begin{subfigure}[t]{0.5\textwidth}
      \centering
      \begin{tikzpicture}
        \standardAxisTwo
        \drawIntervalCC{0}{1}{UiT-blue}
        \drawIntervalOC{0.5}{1.5}{UiT-orange}
        \fill[opacity = 0.2, green,rounded corners=1ex]%
        (0.5,-1.12ex) -- (1, -1.12ex) -- (1, 1.12ex) -- (0.5,1.12ex) -- cycle;
      \end{tikzpicture}
      \caption{$\textcolor{UiT-blue}{U} \cap \textcolor{UiT-orange}{V} = \UiTcorrectColor{[u,1]}$}
      \label{subfig:1.2-4a-MAT-0001-2019-V}
    \end{subfigure}%
    ~ 
    \begin{subfigure}[t]{0.5\textwidth}
      \centering
      \begin{tikzpicture}
        \standardAxisTwo
        \drawIntervalCC{0}{1}{UiT-blue}
        \drawIntervalOC{-0.5}{0.5}{UiT-orange}
        \fill[opacity = 0.2, green,rounded corners=1ex]%
        (0,-1.12ex) -- (0.5, -1.12ex) -- (0.5, 1.12ex) -- (0,1.12ex) -- cycle;
      \end{tikzpicture}
      \caption{$\textcolor{UiT-blue}{U} \cap \textcolor{UiT-orange}{V} = \UiTcorrectColor{[0,w]}$}
      \label{subfig:1.2-4b-MAT-0001-2019-V}
    \end{subfigure}\vspace{\baselineskip}

    \begin{subfigure}[t]{0.5\textwidth}
      \centering
      \begin{tikzpicture}
        \standardAxisTwo
        \drawIntervalCC{0}{1}{UiT-blue}
        \drawIntervalOC{0.25}{0.75}{UiT-orange}
        \fill[opacity = 0.2, green,rounded corners=1ex]%
        (0.25,-1.12ex) -- (0.75, -1.12ex) -- (0.75, 1.12ex) -- (0.25,1.12ex) -- cycle;
      \end{tikzpicture}
      \caption{$\textcolor{UiT-blue}{U} \cap \textcolor{UiT-orange}{V} = \UiTcorrectColor{[u,w]}$}
      \label{subfig:1.2-4c-MAT-0001-2019-V}
    \end{subfigure}%
    ~ 
    \begin{subfigure}[t]{0.5\textwidth}
      \centering
      \begin{tikzpicture}
        \standardAxisTwo
        \drawIntervalCC{0}{1}{UiT-blue}
        \drawIntervalOC{-0.5}{1.5}{UiT-orange}
        \fill[opacity = 0.2, green,rounded corners=1ex]%
        (0,-1.12ex) -- (1, -1.12ex) -- (1, 1.12ex) -- (0,1.12ex) -- cycle;
      \end{tikzpicture}
      \caption{$\textcolor{UiT-blue}{U} \cap \textcolor{UiT-orange}{V} = \UiTcorrectColor{[0,1]}$}
      \label{subfig:1.2-4d-MAT-0001-2019-V}
    \end{subfigure}
    \caption{Illustrasjoner av de fire ulike måtene $U$ og $V$ kan snitte
      hverandre på fra \cref{problem:1.2-4-oving-01-2019-MAT-0001}.}
    \label{fig:1.2-4-MAT-0001-2019-V}
  \end{figure*}
}

\begin{solution}
  Dette er en oppgave som kan virke vanskelig, hvor skal vi starte? Når jeg
  møter en slik oppgave prøver jeg først og forstå hva jeg skal vise. Så først
  lager jeg intervallet $U = [0,1]$. Så er spørsmålet, hvilke ulike måter kan
  jeg snitte $U$ og $V$?  Igjen det kan virke litt vanskelig, men vi spør bare
  hvor vi kan plassere $V$ i forhold til $U$ på tallinjen. Dersom du tenker litt
  ser vi at det \emph{fire muligheter}, disse er illustrert i
  \cref{fig:1.2-4-MAT-0001-2019-V}. La $V = (u,w)$ og anta at $U$ og $V$ har
  punktene $a,b$ til felles, hvor $a \leq b$. Med symboler sier vi at $a,b \in
  U$ og $a,b \in V$. At $a,b$ ligger i $U$ kan vi og skrive som ett intervall
  %
  \begin{equation*}
    0 \leq a \leq b \leq 1,
  \end{equation*}
  %
  og tilsvarende for $V$
  %
  \begin{equation*}
    u \leq a \leq b \leq w.
  \end{equation*}
  %
  Nå vet vi både at $0 \leq a$ og $u \leq a$, men vi vet ikke om $u$ er større
  eller mindre enn null. Så enten har vi $0 \leq u \leq a$ eller $u \leq 0 \leq
  a$.  Tilsvarende så har vi at $b \leq 1$ og $b \leq w$, men vi vet ikke hva
  som er størst av $1$ og $w$, så enten har vi $b \leq 1 \leq w$ eller $b \leq w
  \leq 1$. Dette gir fire muligheter totalt
  %
  \begin{align*}
    0 \UiTcorrectColor{\leq u \leq a \leq b \leq 1}{} \leq w \qquad \text{(se \cref{subfig:1.2-4a-MAT-0001-2019-V})}\\
    u \UiTcorrectColor{\leq 0 \leq a \leq b \leq w}{} \leq 1 \qquad \text{(se \cref{subfig:1.2-4b-MAT-0001-2019-V})}\\
    0 \UiTcorrectColor{\leq u \leq a \leq b \leq w}{} \leq 1 \qquad \text{(se \cref{subfig:1.2-4c-MAT-0001-2019-V})}\\
    u \UiTcorrectColor{\leq 0 \leq a \leq b \leq 1}{} \leq w \qquad \text{(se \cref{subfig:1.2-4d-MAT-0001-2019-V})}
  \end{align*}
  %
  Her ser vi at uansett hvilket av disse fire tilfellene vi har, får vi uansett
  ett intervall.  I denne oppgaven gjorde vi noen antagelser som at vi låste $U$
  til å være $(0,1)$, hvorfor kan vi gjøre dette?
\end{solution}

\newpageLF


%==============================================================================%
%                                 SEKSJON 1.3                                  %
%==============================================================================%

\subsection*{Seksjon 1.3}

%==============================================================================%
%                                PROBLEM 1.3.1                                 %
%==============================================================================%

\begin{problem}[1]
  Forenkle følgende formeluttrykk så mye du kan. Alle bokstaver 
  står for tall.
  \begin{tabenum}[\color{\UiTnumbercolor}\bfseries a\color{black})]%
    \tabenumitem $\displaystyle \frac{x^4 - x^3}{2x + x^2}$
    \hypertarget{subproblem:1.3-1a-MAT-0001-2019-V}{}
    \tabenumitem $\displaystyle \frac{x^3 + 5x^2}{x^3 + 2x}$
    \hypertarget{subproblem:1.3-1b-MAT-0001-2019-V}{}
    \tabenumitem $\displaystyle \frac{x^5 y - x y^2}{xy}$
    \hypertarget{subproblem:1.3-1c-MAT-0001-2019-V}{}
    \\
    \tabenumitem $\displaystyle \frac{x^2 + x}{x + 1}$
    \hypertarget{subproblem:1.3-1d-MAT-0001-2019-V}{}
    \tabenumitem $\displaystyle \frac{1}{x} + \frac{2}{x^2} - \frac{x}{x^3}$
    \hypertarget{subproblem:1.3-1e-MAT-0001-2019-V}{}
    \tabenumitem $\displaystyle x \cdot \frac{5x - 8}{x^2}$
    \hypertarget{subproblem:1.3-1f-MAT-0001-2019-V}{}
    \\
    \tabenumitem $\displaystyle 1 - \frac{x - 1}{x}$
    \hypertarget{subproblem:1.3-1g-MAT-0001-2019-V}{}
    \tabenumitem $\displaystyle \frac{\frac{1}{a}}{\frac{1}{a-b}}$
    \hypertarget{subproblem:1.3-1h-MAT-0001-2019-V}{}
  \end{tabenum}
\end{problem}

\begin{solution}
    Husk å tøye ut før du begynner å gjøre gymnastiske algebraøvelser!
    \begin{flalign*}
        \textit{\hyperlink{subproblem:1.3-1a-MAT-0001-2019-V}{Oppgave~\textcolor{\UiTrefcolor}{1a})}} && 
          &\frac{x^4 - x^3}{2x + x^2}
        = \frac{\textcolor{UiT-orange}{x}\cdot x^3 - \textcolor{UiT-orange}{x} \cdot x^2}{\textcolor{UiT-blue}{x} \cdot 2 + \textcolor{UiT-blue}{x} \cdot x}
        = \frac{\textcolor{UiT-orange}{x}}{\textcolor{UiT-blue}{x}} \cdot\frac{x^3 - x^2}{2 + x}
        = \frac{x^2(x - 1)}{x + 2} \\
%
        \textit{\hyperlink{subproblem:1.3-1b-MAT-0001-2019-V}{Oppgave~\textcolor{\UiTrefcolor}{1b})}} && 
        &\frac{x^3 + 5x^2}{x^3 + 2x}
        = \frac{\textcolor{UiT-orange}{x} \cdot x^2 + \textcolor{UiT-orange}{x} \cdot 5x}{\textcolor{UiT-blue}{x} \cdot x^2 + \textcolor{UiT-blue}{x} \cdot 2}
        = \frac{\textcolor{UiT-orange}{x}}{\textcolor{UiT-blue}{x}} \cdot \frac{x^2 + 5x}{x^2 + 2}
        = \frac{x^2 + 5x}{x^2 + 2} \\
%
        \textit{\hyperlink{subproblem:1.3-1c-MAT-0001-2019-V}{Oppgave~\textcolor{\UiTrefcolor}{1c})}} && 
        &\frac{x^5y - x y^2}{xy}
        = \frac{\textcolor{UiT-orange}{xy} \cdot x^4 - \textcolor{UiT-orange}{y}\cdot y}{\textcolor{UiT-blue}{xy}}
        = \frac{\textcolor{UiT-orange}{xy}}{\textcolor{UiT-blue}{xy}} (x^4 -y) 
        = x^4 - xy \\
%
        \textit{\hyperlink{subproblem:1.3-1d-MAT-0001-2019-V}{Oppgave~\textcolor{\UiTrefcolor}{1d})}} && 
           &\frac{x^2 + x}{x + 1}
         = \frac{\textcolor{UiT-orange}{x} \cdot x + \textcolor{UiT-orange}{x}}{x + 1}
         = \frac{\textcolor{UiT-orange}{x} (\textcolor{UiT-blue}{x + 1})}{\textcolor{purple}{x + 1}}
         = x \cdot \frac{\textcolor{UiT-blue}{x + 1}}{\textcolor{purple}{x + 1}}
         = x \\
%
          \textit{\hyperlink{subproblem:1.3-1e-MAT-0001-2019-V}{Oppgave~\textcolor{\UiTrefcolor}{1e})}} && 
          &\frac{1}{x} + \frac{2}{x^2} - \frac{x}{x^3}
          = \frac{x}{x}\frac{1}{x} + \frac{2}{x^2} - \frac{1}{x^2}
          = \frac{x + 2 - 1}{x^2}
          = \frac{x + 1}{x^2} \\
%
          \textit{\hyperlink{subproblem:1.3-1f-MAT-0001-2019-V}{Oppgave~\textcolor{\UiTrefcolor}{1f})}} && 
          &x \cdot \frac{5x - 8}{x^2}
          = \textcolor{UiT-orange}{x} \cdot \frac{5x - 8}{\textcolor{UiT-blue}{x} \cdot x}
          = \frac{\textcolor{UiT-orange}{x}}{\textcolor{UiT-blue}{x}} \frac{5x - 8}{x}
          = 5 - \frac{8}{x} \\
%
          \textit{\hyperlink{subproblem:1.3-1g-MAT-0001-2019-V}{Oppgave~\textcolor{\UiTrefcolor}{1g})}} && 
            &1 - \frac{x-1}{x}
          = 1 - \biggl( \frac{x}{x} - \frac{1}{x}\biggr)
          = 1 - 1 + \frac{1}{x}
          = \frac{1}{x} \\
%
          && 
            &1 - \frac{x-1}{x}
          = \frac{x}{x} - \frac{x-1}{x}
          = \frac{x -(x-1)}{x}
          = \frac{x - x + 1}{x}
          = \frac{1}{x} \\
%
           \textit{\hyperlink{subproblem:1.3-1h-MAT-0001-2019-V}{Oppgave~\textcolor{\UiTrefcolor}{1h})}} && 
            &\frac{\frac{1}{a}}{\frac{1}{a-b}}
            = \frac{\frac{1}{a}}{\frac{1}{a-b}} \frac{a-b}{a-b}
            = \frac{\frac{1}{a} \cdot (a-b)}{\frac{1}{a-b} \cdot (a-b)}
            = \frac{\frac{a-b}{a}}{1}
            = \frac{a}{a}-\frac{b}{a} = 1 - \frac{b}{a}
    \end{flalign*}
\end{solution}

%==============================================================================%
%                                PROBLEM 1.3.2                                 %
%==============================================================================%

\begin{problem}
  Et stoff har massetetthet $\SI{12}{\g\per\cm\cubed}$, det vil si $12$~gram per
  kubikkcentimeter. Hvor mye veier $\SI{10}{\m\cubed}$ av dette stoffet?
\end{problem}

\begin{solution}
  La oss først bytte til enheter til. Siden $\SI{1}{\cm} = \SI{e-2}{\m}$ får vi
  %
  \begin{equation*}
      \frac{\SI{12}{\g}}{\si{\cm\cubed}}
    = \frac{\SI{12}{\g}}{(\SI{e-2}{\m})^3}
    = \frac{\SI{12}{\g}}{\SI{e-6}{\m\cubed}}
    = \SI{12e6}{\g\per\m\cubed} 
    = \SI{12e7}{\g}/\SI{10}{\m\cubed} 
  \end{equation*}
  %
  Altså veier $\SI{10}{\m\cubed}$ av dette stoffet \UiTcorrectColor{\SI{1.2e8}{\g}}
  eller \UiTcorrectColor{\SI{1.2e5}{\kg}} eller \UiTcorrectColor{\SI{120}{\tonne}}.
\end{solution}

\newpageLF

%==============================================================================%
%                                PROBLEM 1.3.3                                 %
%==============================================================================%

\begin{problem}
  Hvor mange siffer har
  %
  \begin{subproblem}{3}
    \item $\num{-14.5400}$ \UiTcorrect{$6$}
      \label{subproblem:1.3-3a-MAT-0001-2019-V}
    \item $\num{-0.14e3}$ \UiTcorrect{$2$}
      \label{subproblem:1.3-3b-MAT-0001-2019-V}
    \item $\num{0.004}$ \UiTcorrect{$1$}
      \label{subproblem:1.3-3c-MAT-0001-2019-V}
  \end{subproblem}
\end{problem}

\begin{solution}
  Når vi skal regne gjeldende siffer så teller vi alltid fra og med det første
  tallet som ikke er null, og til og med den siste nullen etter komma. Dette er
  fordi det er stor forskjell på $\num{120,0}$ og $\num{120,000000}$ i siste
  tilfellet har vi gjort langt, langt mer nøyaktige målinger og har en mye
  lavere usikkerhet. \medskip

  \Opg{subproblem:1.3-3a-MAT-0001-2019-V} Her har vi to siffer før
  komma, og fire etter, altså \UiTcorrectColor{$6$ siffer} totalt.
  \Opg{subproblem:1.3-3b-MAT-0001-2019-V} Her har vi null siffer før
  komma, og to etter, altså \UiTcorrectColor{$2$ siffer} totalt.
  \Opg{subproblem:1.3-3c-MAT-0001-2019-V} Her har vi null siffer før
  komma, og ett etter, altså \UiTcorrectColor{$1$ siffer} totalt
\end{solution}

%==============================================================================%
%                                PROBLEM 1.3.4                                 %
%==============================================================================%

\begin{problem}
  Sidene i et rektangulært svømmebasseng er målt i henholdsvis $\num{5.3}$~meter
  og $\num{2.997}$~meter med usikkerhet. Finn arealet av bassengets overflate,
  fornuftig avrundet. \UiTcorrect{$16$~meter}
\end{problem}

\begin{solution}
  Når vi skal regne med enheter er en fin huskeregel at vi alltid skal ta med
  like mange gjeldende siffer som det tallet med \emph{færrest} gjeldende
  siffer. I dette tilfellet skal vi altså ha med \emph{to} gjeldende siffer.

  Merk at når vi gjør selve utregningen tar vi med alle sifferene, og gjør
  avrundingen til sist slik at vi unngår avrundingsfeil.
  %
  \begin{equation*}
    A = 2b + 2l = 2(\SI{5.3}{\m} + \SI{2.997}) = \SI{15.8841}{\m}
  \end{equation*}
  %
  Arealet av omkretsen til bassenget er altså $\num{16}$~meter. Hvor
  vi da rundet opp $\num{15.8}$ til $\num{16}$. Husk at dersom vi ønsker å runde
  av til $k$ gjeldende siffer, så ser vi bare på de $k+1$ første sifferene i
  tallet og runder av disse.
\end{solution}

%==============================================================================%
%                                PROBLEM 1.3.5                                 %
%==============================================================================%

\begin{problem}
  Du har fått tallet $\num{871.5}$ som resultat av en måling, men vurdering av
  usikkerheten tilsier at du bare har grunnlag for å oppgi svaret med to
  gjeldende siffer. Hvordan bør svaret skrives?\UiTcorrect[~]{870}[~]Hva om du bare 
  kan bruke ett gjeldende siffer? \UiTcorrect{900}
\end{problem}

\begin{solution}
  Dersom vi skal skrive svaret med to gjeldende siffer så er det tallet $87$ vi
  skal bruke, siden vi runder ned $1$. Så tallet kan skrives som
  $\num{8.7e2}$ eller $\num{870}$.

  Dersom svaret skal skrives med ett gjeldende siffer så runder vi opp $87$ til
  $90$. Så tallet kan skrives som $\num{9e2}$ eller
  $\num{900}$.
\end{solution}

\newpageLF

%==============================================================================%
%                       BLANDEDE OPPGAVER TIL KAPITTEL 1                       %
%==============================================================================%

\subsection*{Blandede oppgaver til kapittel 1}

%==============================================================================%
%                               PROBLEM 1.X.14                                 %
%==============================================================================%

\begin{problem}[14]
  \label{problem:1-mix-14-MAT-0001-2019-V} Anta $a < -3$. Finn $\abs{a}$ uttrykt
  ved $a$. \UiTcorrect{$\abs{a} = -a$}
\end{problem}

\begin{solution}
  \Opg{problem:1-mix-14-MAT-0001-2019-V}: Det enkleste blir nok å
  huske på definisjonen av absoluttverdien
  %
  \begin{equation*}
    \label{eq:1-mix-14-MAT-0001-2019-V-def-absolute-value}
    \abs{x} = 
    \begin{cases}
      -x & \ \text{hvis} \ x \leq 0, \\
      \phantom{-} x & \ \text{hvis} \ x > 0
    \end{cases}.
  \end{equation*}
  %
  Så siden $a < -3$ har vi per definisjon at $\abs{a} = -a$. Vi kan
  og gjøre dette med ett eksempel, husk at absoluttverdien av ett tall alltid
  skal være \emph{positivt}. La $a = -5$ da er det klart at $a < -3$. Så regner
  vi ut$\abs{-5} = -(-5) = 5$ så her ser vi at absoluttverdien git ut et
  positivt svar som var det vi ønsket! Her brukte vi bare det vi fant ut altså
  at $\abs{x} = -x$ for negative $x$.
\end{solution}

%==============================================================================%
%                               PROBLEM 1.X.15                                 %
%==============================================================================%

\begin{problem}
  For hvilke verdier av $x$ har vi
    \begin{subproblem}{3}
      \item $x < \abs{x}$ \UiTcorrect{$x<0$}
        \label{subproblem:1-mix-16a-MAT-0001-2019-V}
      \item $\abs{x} \leq x$ \UiTcorrect{$x\geq0$}
        \label{subproblem:1-mix-16b-MAT-0001-2019-V}
      \item $\abs{x} < x$ \UiTcorrect{$\varnothing$}
        \label{subproblem:1-mix-16c-MAT-0001-2019-V}
    \end{subproblem}
\end{problem}

\hvisLF{%
  \begin{figure}[htbp!]
    \centering
    \begin{tikzpicture}
      \begin{axis}[
        UiTplotstyle,
        ytick={-6,-4,...,6},
        xtick={-6,-4,...,6},
        % x tick label style={below right},
        % y tick label style={above left,},
        ymin=-4.5,
        ymax=4.5,
        xmin=-4.5,
        xmax=4.5,
        domain=-10:10,
        after end axis/.code={
          \path (axis cs:0,0) 
          node [anchor=north west,yshift=-0.075cm] {0}
          node [anchor=south east,xshift=-0.075cm] {0};
        }
        ]
        \addplot+[name path=F, domain=-5:0, mark=none, white] {-x};
        \addplot+[name path=G, domain= 0:5, mark=none, white] { x};

        \addplot+[name path=H, ultra thick,mark=none, UiT-orange] {x};

        \addplot+[domain=-5:0, only marks,mark=star,mark options={color=UiT-blue,fill=white,scale=1}] {-x};
        \addplot+[domain= 0:5, only marks,mark=star,mark options={color=UiT-blue,fill=white,scale=1}] { x};

        \addplot[opacity=0.2,green]fill between[of=F and H, soft clip={domain=-5:0}];
      \end{axis}
    \end{tikzpicture}
    \caption{$\textcolor{UiT-blue}{\abs{x}}$ og $\textcolor{UiT-orange}{x}$}
    \label{subfig:1-mix-16-MAT-0001-2019-V}. 
  \end{figure}
}

\begin{solution}
  Det enkleste blir nok å tegne $x$ og $\abs{x}$ også drøfte hvilken som
  er størst på hvilket intervall. En enkel skisse er vist i
  \cref{subfig:1-mix-16-MAT-0001-2019-V}. \medskip

  \Opg{subproblem:1-mix-16a-MAT-0001-2019-V} Fra figuren så
  viser det grønne området viser hvor $x<\abs{x}$. Dette skjer når $x<0$,
  men la oss vise det ved å bruke definisjonen av $\abs{x}$ 
  %
  \begin{equation*}
    \abs{x}>x \ \Leftrightarrow \ 
    \begin{cases}
      -x>x & \text{hvis} \ x \leq 0, \\
      \phantom{-} x>x & \text{hvis} \ x > 0
    \end{cases}.
  \end{equation*}        
  %
  Hvor vi ser at likningen har ingen løsninger når $x>0$. I det andre
  tilfellet har vi $-x > x \ \Rightarrow \ 0 > 2x \ \Rightarrow 0 > x$ som
  er det samme som at $x < 0$.
  \medskip

  \Opg{subproblem:1-mix-16b-MAT-0001-2019-V} Igjen fra
  \cref{subfig:1-mix-16-MAT-0001-2019-V} ser vi at eneste gangen hvor $x$
  er større eller lik $\abs{x}$ er når $x\geq 0$. Dette gir mye mening da
  $\abs{x} = x$ når $x$ er positiv. 
  %
  \begin{equation*}
    \abs{x} \leq x \ \Leftrightarrow \ 
    \begin{cases}
      -x \leq x & \text{hvis} \ x \leq 0, \\
      \phantom{-} x \leq x & \text{hvis} \ x > 0
    \end{cases}
  \end{equation*}        
  %
  Hvor første likning $-x \leq x$ er gyldig for $x = 0$, mens $x \leq x$
  er gyldig for alle $x>0$. Dette gir til sammen den samme løsningen som
  før, altså at $\abs{x} \leq x \ \Rightarrow \ x \geq 0$.
  \medskip

  \Opg{subproblem:1-mix-16c-MAT-0001-2019-V} Fra
  \cref{subfig:1-mix-16-MAT-0001-2019-V} ser vi at likningen $\abs{x} < x$
  ikke kan ha noen løsninger. Dette er fordi funksjonen $x$ aldri har
  høyere verdi enn $\abs{x}$ bare like stor eller lavere. Ved å bruke
  definisjonen får vi helt tilsvarende
  %
  \begin{equation*}
    \abs{x} < x \ \Leftrightarrow \ 
    \begin{cases}
      -x < x & \text{hvis} \ x \leq 0, \\
      \phantom{-} x < x & \text{hvis} \ x > 0
    \end{cases}
    \ \Leftrightarrow \ 
    \begin{cases}
      x > 0 & \text{hvis} \ x \leq 0, \\
      0 < 0 & \text{hvis} \ x > 0
    \end{cases}.
  \end{equation*}        
  %
  Hvor vi ser at begge tilfellene er umulige og oppnå slik at det er
  ingen verdier for $x$ slik at $\abs{x} < x$.
\end{solution}

%==============================================================================%
%                               PROBLEM 1.X.17                                 %
%==============================================================================%

\begin{problem}[17]
  Skriv mengdene som en union av to intervaller:
  \begin{multicols}{2}
    \begin{subproblem}
      \label{subproblem:1-mix-17a-MAT-0001-2019-V}
      $\Set{x \in \R \given \abs{x} > 2}$
    \end{subproblem}
    \begin{subproblem}
      \label{subproblem:1-mix-17b-MAT-0001-2019-V}
      $\Set{x \in \R \given \abs{x - 3} > 2}$
    \end{subproblem}
  \end{multicols}
\end{problem}

\hvisLF{%
  \begin{figure*}[htbp!]
    \centering
    \begin{subfigure}[t]{0.45\textwidth}
      \centering
      \begin{tikzpicture}
        \begin{axis}[
          UiTplotstyle,
          width=1.2\textwidth,
          ytick={-6,-4,...,6},
          xtick={-6,-4,...,6},
          ymin=-0.5,
          ymax=4.5,
          xmin=-4.5,
          xmax=4.5,
          domain=-10:10,
          after end axis/.code={
            \path (axis cs:0,0) 
%            node [anchor=north west,yshift=-0.075cm] {0};
            node [anchor=south east,xshift=-0.075cm] {0};
          }
          ]
          \addplot+[name path=F,domain=-5:0, ultra thick, mark=none,UiT-blue] {-x}
          node [pos=0.6, above] {$\abs{x}$};
          \addplot+[name path=G,domain= 0:5, ultra thick, mark=none,UiT-blue] {x};
          \addplot+[name path=H,ultra thick, mark=none, UiT-orange] {2}
          node [pos=0.65, below] {$y = 2$};
          \addplot[opacity=0.2,green]fill between[of=F and H, soft clip={domain=-5:-2}];
          \addplot[opacity=0.2,green]fill between[of=G and H, soft clip={domain=2:5}];
        \end{axis}
      \end{tikzpicture}
      \caption{\UiTcorrectColor{$(-\infty,-2)\cup(2,\infty)$}{}}
      \label{subfig:1-mix-17a-MAT-0001-2019-V}
    \end{subfigure}
    %
    \begin{subfigure}[t]{0.45\textwidth}
      \centering
      \begin{tikzpicture}
        \begin{axis}[
          UiTplotstyle,
          width=1.2\textwidth,
          ytick={-6,-4,...,6},
          xtick={0,1,3,5,7},
          ymin=-0.5,
          ymax=4.5,
          xmin=-1.5,
          xmax=7.5,
          domain=-10:10,
          after end axis/.code={
            \path (axis cs:0,0) 
%            node [anchor=north west,yshift=-0.075cm] {0}
            node [anchor=south east,xshift=-0.075cm] {0};
          }
          ]
          \addplot+[name path=F,domain=-2:3, ultra thick, mark=none,UiT-blue] {3-x}
          node [pos=0.75, above,yshift=1em] {$\abs{x-3}$};
          \addplot+[name path=G,domain= 3:8, ultra thick, mark=none,UiT-blue] {x-3};
          \addplot+[name path=H,ultra thick, mark=none, UiT-orange] {2}
          node [pos=0.8, below] {$y = 2$};
          \addplot[opacity=0.2,green]fill between[of=F and H, soft clip={domain=-2:1}];
          \addplot[opacity=0.2,green]fill between[of=G and H, soft clip={domain=5:8}];
        \end{axis}
      \end{tikzpicture}
      \caption{\UiTcorrectColor{$(-\infty,1)\cup(5,\infty)$}{}}
      \label{subfig:1-mix-17b-MAT-0001-2019-V}
    \end{subfigure}
    \caption{}
    \label{fig:1-mix-17-MAT-0001-2019-V}
  \end{figure*}
}

\begin{solution}
  Her ønsker vi å skrive om mengdene som en union av to intervall.  En smart idè
  kan være å tegne funksjonene som vi har gjort i
  \cref{fig:1-mix-17-MAT-0001-2019-V}. Vi ønsker området hvor $\abs{x}$ ligger
  over $y = 2$, disse to intervallene er da markert med grønn i
  \cref{subfig:1-mix-17a-MAT-0001-2019-V,subfig:1-mix-17b-MAT-0001-2019-V}.
  Alternativt kan intervallene også bestemmes via regning: \medskip

  \Opg{subproblem:1-mix-17a-MAT-0001-2019-V} Den vanligste metoden er
  å bruke definisjonen av absoluttverdien
  %
  \begin{equation*}
    \abs{x} > 2
    \ \Leftrightarrow \ 
    \begin{cases}
      -  x > 2 \ \text{hvis} \ x < 0 \\
      \phantom{-} x > 2 \ \text{hvis} \ x \geq 0
    \end{cases}
    \ \Leftrightarrow \ 
    \begin{cases}
      x < - 2 \ \text{hvis} \ x < 0 \\
      x > \phantom{-} 2 \ \text{hvis} \ x \geq 0
    \end{cases}.
  \end{equation*}
  %
  Hvor vi kan lese ut at intervallene våre er $(-\infty,-2) \cup (2,\infty)$.
  Husk at når vi ganger med $-1$ så snur vi alle ulikheter.  Altså at 
  %
  \begin{equation*}
    -x > 2 \ \Rightarrow \ x < -2
  \end{equation*}
  %
  Dette kan vi og se dersom vi \bquote{flytter\footnote{Eller trekker fra det
  samme på begge sider som de mer korrekte matematikkerene ville sagt det.}}
  over $x$, og deretter $2$
  %
  \begin{equation}
    \label{eq:1-mix-17b-MAT-0001-2019-V-ulikhet}
    -x > 2 
    \ \Rightarrow \  0 > x + 2 
    \ \Rightarrow \ -2 > x 
    \ \Rightarrow \  x < 2.
  \end{equation}
  %
  \Opg{subproblem:1-mix-17b-MAT-0001-2019-V} Igjen fra definisjonen
  har vi
  %
  \begin{equation*}
    \abs{x-3} > 2
    \ \Leftrightarrow \ 
    \begin{cases}
      -(x-3) > 2 \ \text{hvis} \ x < 0 \\
      \phantom{-}(x-3) > 2 \ \text{hvis} \ x \geq 0
    \end{cases}
    \Leftrightarrow \ 
    \begin{cases}
      x < 1\ \text{hvis} \ x < 0 \\
      x > 5 \ \text{hvis} \ x \geq 0
    \end{cases}.
  \end{equation*}
  %
  Hvor vi kan lese ut at intervallene våre er $(-\infty,1)\cup(5,\infty)$.  Å
  vise at $-(x-3) > 2 \ \Rightarrow \ x < 1$ kan vises enten ved å først snu
  ulikheten ved å gange med $-1$, eller gjennomføre de samme stegene som i
  \cref{eq:1-mix-17b-MAT-0001-2019-V-ulikhet}; men dette overlates som en artig
  øvelse til leseren ;-)
\end{solution}

%==============================================================================%
%                               PROBLEM 1.X.20                                 %
%==============================================================================%

\begin{problem}[20]
  Avgjør om følgende mengder er like.
  \begin{subproblem}
    \label{subproblem:1-mix-20a-MAT-0001-2019-V}
    $A = \Set{\Set{1,2}}$ og $B = \Set{1, 2}$ 
    \UiTwrong{$A \neq B$}
  \end{subproblem}
  \begin{subproblem}
    \label{subproblem:1-mix-20b-MAT-0001-2019-V}
    $A = \Set{1, 2, 1, 1, 4}$ og $B = \Set{4, 2, 1}$
    \UiTcorrect{$A = B$}
  \end{subproblem}
  \begin{subproblem}
    \label{subproblem:1-mix-20c-MAT-0001-2019-V}
    $A = \Set{1, 2, 3, 4}$ og $B = \Set{4, 3, 2, 1}$
    \UiTcorrect{$A = B$}
  \end{subproblem}
\end{problem}

\begin{solution}
  Vi begynner med å skrive opp definisjonen for at to mengder skal være like
  %
  \begin{equation*}
    A = B 
    \quad \Leftrightarrow \quad 
    A \subseteq B \ \text{og} \ B \subseteq A
  \end{equation*}
  %
  med andre ord så må \emph{alle} elementene i $A$ være inneholdt i $B$, og
  \emph{alle} elementene i $B$ må være inneholdt i $A$. \medskip

  \Opg{subproblem:1-mix-20a-MAT-0001-2019-V} Vi kan se med en gang at
  elementet $\Set{1,2} \in A$ ikke ligger i $B$ ($\Set{1,2} \not\in B$) slik at
  $A \neq B$.  \medskip

  \Opg{subproblem:1-mix-20b-MAT-0001-2019-V} I en mengde så er alle
  elementene unike, ved å gjøre dette på $A$, ser vi at $A = \Set{4,2,1}$ så $A$
  og $B$ er identiske mengder.  Slik at $A = B$.  \medskip

  \Opg{subproblem:1-mix-20c-MAT-0001-2019-V} Her kan en enten si at
  \UiTcorrectColor{$A = B$} direkte siden rekkefølgen av elementer i en mengde er
  irrelevant, eller sjekke at alle elementene i $A$ også ligger i $B$, og at
  alle elementene i $B$ også ligger i $A$.
\end{solution}

\newpageLF

%==============================================================================%
%                               PROBLEM 1.X.21                                 %
%==============================================================================%

\begin{problem}
  La $A = \Set{4, 5, 6, 7, 8, 9}$ og $B = \Set{2, 4, 8, 10}$. Finn $A \cap B$
  og $A \cup B$.
\end{problem}

\begin{solution}
  \textcolor{UiT-main}{$A \cap B$}: Ønsker å finne \emph{snittet} mellom $A$ og
  $B$ altså bare de som ligger i begge mengdene. Med andre ord bare det de har
  til felles.
  %
  \begin{align*}
    A \cap B &= \Set{4, 8}.
    %
    \intertext{\textcolor{UiT-main}{$A \cup B$}: Ønsker å finne \emph{unionen}
    mellom $A$ og $B$ altså alt som ligger i begge mengdene til sammen.}
    %
    A \cup B &= \Set{2,4, 5, 6, 7, 8, 9, 10}.
  \end{align*}
\end{solution}

%==============================================================================%
%                               PROBLEM 1.X.22                                 %
%==============================================================================%

\begin{problem}
  Avgjør om $A$ er inneholdt i $B$.
  \begin{subproblem}
    \label{subproblem:1-mix-22a-MAT-0001-2019-V}
    $A = \Set{2, 3, 5}$ og $B = \Set{2, 3, 4, 5}$
    \UiTcorrect{$A \subseteq B$}
  \end{subproblem}
  \begin{subproblem}
    \label{subproblem:1-mix-22b-MAT-0001-2019-V}
    $A = \Set{2, 3, 5}$ og $B = \Set{2, 3, 5}$
    \UiTcorrect{$A \subseteq B$}
  \end{subproblem}
  \begin{subproblem}
    \label{subproblem:1-mix-22c-MAT-0001-2019-V}
    $A = \Set{1, 2}$ og $B = \Set{\Set{1, 2}}$
    \UiTwrong{$A \not\subseteq B$}
  \end{subproblem}
  \begin{subproblem}
    \label{subproblem:1-mix-22d-MAT-0001-2019-V}
    $A = \varnothing$ og $B = \Set{2, 3, 4, 5}$
    \UiTcorrect{$A \subseteq B$}
  \end{subproblem}
  \begin{subproblem}
    \label{subproblem:1-mix-22e-MAT-0001-2019-V}
    $A = \varnothing$ og $B = \varnothing$
    \UiTcorrect{$A \subseteq B$}
  \end{subproblem}
\end{problem}

\begin{solution}
  For å finne ut om mengde $A$ er inneholdt i $B$ så må alle $x \in A$ også
  ligge i $B$. \medskip

  \Opg{subproblem:1-mix-22a-MAT-0001-2019-V} Her
  må vi sjekke om $2$, $3$ og $5$ ligger i $B$
  %
  \begin{equation*}
    A = \Set{\UiTcorrectColor{2}, \UiTcorrectColor{3}, \UiTcorrectColor{5}} 
    \ \text{og} \ 
    B = \Set{\UiTcorrectColor{2}, \UiTcorrectColor{3}, 4, \UiTcorrectColor{5}} \
    \text{slik at} \ A \subseteq B.
  \end{equation*}
  %
  \Opg{subproblem:1-mix-22b-MAT-0001-2019-V} Her må vi sjekke om $2$,
  $3$ og $5$ ligger i $B$
  %
  \begin{equation*}
    A = \Set{\UiTcorrectColor{2}, \UiTcorrectColor{3}, \UiTcorrectColor{5}} 
    \ \text{og} \ 
    B = \Set{\UiTcorrectColor{2}, \UiTcorrectColor{3}, \UiTcorrectColor{5}}  \
    \text{slik at} \ A \subseteq B.
  \end{equation*}
  %
  \Opg{subproblem:1-mix-22c-MAT-0001-2019-V} Denne kan være litt
  tricky for her må vi sjekke om $B$ inneholder $1$ og $2$, men det gjør ikke
  $B$.  Det eneste $B$ inneholder er mengden $\Set{1, 2}$ som hverken er lik $1$
  eller $2$!
  %
  \begin{equation*}
    A = \Set{\UiTcorrectColor{1}, \UiTcorrectColor{2}} 
    \ \text{og} \ 
    B = \Set{\UiTwrongColor{\Set{1, 2}}} \
    \text{slik at} \ A \not\subseteq B.
  \end{equation*}
  %
  \Opg{subproblem:1-mix-22d-MAT-0001-2019-V} Vi ønsker at alle
  elementene i $A$ skal være inneholdt i $B$. Siden $A$ ikke inneholder noen
  elementer (det er den tomme mengden), så inneholder $B$ alle elementene til
  $A$ (som er ingen).
  %
  \begin{equation*}
    A = \varnothing
    \ \text{og} \ 
    B = \Set{2, 3, 4, 5} \
    \text{slik at} \ A \subseteq B.
  \end{equation*}
  %
  \Opg{subproblem:1-mix-22e-MAT-0001-2019-V} Her kan vi bruke akkuratt
  samme argument som i svaret på forrige deloppgave, men la oss i stedet
  presentere ett alternativt argument.  Snittet av to mengder er en undermengde
  til de originale mengdene
  %
  \begin{equation*}
    A \cap B = C \Rightarrow C \subseteq A \ \text{og} \ C \subseteq B.
  \end{equation*}
  %
  Videre har vi at
  $A \cap \varnothing = \varnothing$. Dette betyr at $\varnothing$ er en
  undermengde av $A$ ($\varnothing \subseteq A)$ og at $\varnothing$ er en
  undermengde av $\varnothing$ ($\varnothing \subseteq \varnothing$). Som var
  akkurat det oppgaven spurte om
  %
  \begin{equation*}
    A = \varnothing
    \ \text{og} \ 
    B = \varnothing \
    \text{slik at} \ A \subseteq B.
  \end{equation*}
\end{solution}

\end{document}
